%%
%% fall-lecture04.tex
%% 
%% Made by Alex Nelson <pqnelson@gmail.com>
%% Login   <alex@lisp>
%% 
%% Started on  2026-10-08T09:56:02-0700
%% Last update 2026-10-08T09:56:02-0700
%% 

\lecture{}

\begin{node}
First-countability is local: every \emph{point} has a countable
neighborhood base. Second-countability is global: the entire space has
a countable base. We saw how metrizability implies $(A1)$, we also saw
that $(A2)$ implies $(A1)$.
\end{node}

\begin{definition}
The topological space $(X,\topology)$ is \define{Separable} if there
exists a countable dense set.
\end{definition}

\begin{definition}
We say a topological space $(X,\topology)$ is \define{Lindel\"{o}f} if
every open cover contains a countable subcover.
\end{definition}

\begin{example}
The Lebesgue space $\ell^{\infty}$ is $(A1)$ since it is a metric
space, but $\ell^{\infty}$ is not $(A2)$.
\end{example}

\begin{example}
The space $(\RR,\topology)$ with the topology generated by half-open
intervals $(a,b]$. Then $\topology$ is first-countable $(A1)$ but it
is not second-countable $(A2)$. Further, $\topology$ is separable
(e.g., the rational numbers are dense in this topology). But
$\topology$ is not metrizable (because, for metric spaces, $(A2)$ is
equivalent to separability).
\end{example}

\begin{theorem}
Second-countable spaces are separable.
\end{theorem}

\begin{proof}
Let $\{B_{j}\}$ be a countable base. Then we choose a point $x_{j}\in B_{j}$
for each $j$, then we claim $\{x_{j}\}$ is dense and therefore the
topology is separable.

Assume for contradiction $\{x_{j}\}$ is not dense. Then there exists a
$U\subset X$ open subset such that $U\cap\{x_{j}\}=\emptyset$ (i.e.,
$U$ does not contain $x_{j}$ for any $j$). But since $\{B_{j}\}$ is a
base, there must be some $B_{k}\subset U$ which means $x_{k}\in U$,
hence contradiction.
\end{proof}

\begin{theorem}
A metric space is second-countable if and only if it is separable.
\end{theorem}
(Suffices to prove: separability implies second-countable.)

\begin{proof}
Let $(X,\topology)$ be metrizable. Assume it is separable. We want to
show $\topology$ is second-countable. Let $\{x_{k}\}$ be a countable
dense set. Then we claim $\{B(x_{k},1/k)\}$ is a base for
$(X,\tau)$. This is obviously a countable set.
\end{proof}

\begin{theorem}
Second-countable spaces are Lindel\"{o}f.
\end{theorem}

\begin{proof}
Let $(X,\topology)$ be second-countable. Let $\mathcal{B}=\{B_{j}\}$ be a
countable base for the topology. Let $\mathcal{F}$ to be any open
cover. (We want to find a countable open subcover.)

Consider
\begin{equation}
\mathcal{B}_{\mathcal{F}}=\{B\in\mathcal{B}\mid\exists A\in\mathcal{F}\ldotp B\subset A\}.
\end{equation}
For each $j$, choose $A_{j}\in\mathcal{F}$ such that $B_{j}\subset A_{j}$.
Then these $\{A_{j}\}$ form a subcover of $\mathcal{F}$, and it is
countable, hence the claim.
\end{proof}

\subsection{Separation}

\begin{definition}
Let $(X,\topology)$ be a topological space.
\begin{enumerate}
\item We say it is $T_{1}$ if: for all $x,y\in X$ such that $x\neq y$,
  there exists a $U\in\topology$ such that $x\notin U$ and $y\in U$.
  (Equivalently, every singleton set is closed [``all points are closed''])
\item We say it is $T_{2}$ if: for all $x,y\in X$ such that $x\neq y$,
  there exist disjoint neighborhoods $U\ni x$ and $V\ni y$ such that
  $U\cap V=\emptyset$. We also say the space is \define{Hausdorff} in
  this situation.
\item We say it is $T_{3}$ if: it is $T_{1}$ and for any $x\in X$ and
  any closed subset $C\subset X$ such that $x\notin C$, there exists
  disjoint neighborhoods $x\in U$ and $C\subset V$ such that $U\cap V=\emptyset$.
  We also call such spaces \define{Regular}.
\item We say it is $T_{4}$ if: it is $T_{1}$ and for any disjoint
  closed subsets $F_{1}\subset X$ and $F_{2}\subset X$ there exists
  disjoint neighborhoods $U_{1}\supset F_{1}$ and $U_{2}\supset F_{2}$
  such that $U_{1}\cap U_{2}=\emptyset$. In this case, we also call
  such spaces \define{Normal}.
\end{enumerate}
\end{definition}

\begin{node}
Observe that $T_{4}\implies T_{3}\implies T_{2}\implies T_{1}$.
For us, the most important axioms are $T_{2}$ (Hausdorff spaces) and
$T_{4}$ (normal spaces). A space is Hausdorff if and only if every net
has at most one limit. And normal spacecs guarantees the existence of
many continuous functions ($T_{3}$ is not quite enough, the $T_{4}$
setting is far more natural).
\end{node}

\begin{lemma}
Let $(X,\topology)$ be a metrizable space. Let $\rho$ be a metric on
$X$ generating the topology $\topology$. Let $C\subset X$ be a closed
subset. For any $x\in X$, $\rho(x,C)=0$ implies $x\in C$.
\end{lemma}

\begin{lemma}
Let $(X,\topology)$ be a metrizable space. Let $\rho$ be a metric on
$X$ generating the topology $\topology$. Let $C\subset X$ be a closed
subset. The function $\rho(-,C)$ is continuous.
\end{lemma}

\begin{proof}[Proof sketch]
The triangle inequality $|\rho(x,C)-\rho(y,C)|\leq\rho(x,y)$ gives the result.
\end{proof}

\begin{theorem}
All metrizable spaces are normal.
\end{theorem}

\begin{proof}
Let $A, B\subset X$ be disjoint closed subsets. Then we define the function
\begin{equation}
f(x) = \frac{\rho(x,A)-\rho(x,B)}{\rho(x,A)+\rho(x,B)}.
\end{equation}
Observe the denominator is never zero (by the first lemma). We see
$f=1$ on $B$ and $f=-1$ on $A$ and $f$ is continuous. Then consider
open disjoint sets $f^{-1}((-\infty,0))$ and $f^{-1}((0,\infty))$
which contain $A$ and $B$ respectively. The result follows.
\end{proof}

\begin{remark}
The main result will be Urysohn's lemma, which is used in proving
Urysohn's metrizability theorem. But first, let us look at some examples.
\end{remark}

\begin{example}
The trivial topology is not $T_{1}$.
\end{example}

\begin{example}
The cofinite topology is $T_{1}$ but it is not Hausdorff, which
guarantees $T_{1}\neq T_{2}$.
\end{example}

\begin{example}
Consider the product topology onthe function space
$[0,1]^{[0,1]}$. This is Hausdorff $(T_{2})$---more generally, the
product of Hausdorff spaces is Hausdorff.
\end{example}

\begin{proposition}
Let $(X_{\lambda},\topology_{\lambda})_{\lambda\in\Lambda}$ be a
family of topological spaces. Then the product space $\prod X_{\lambda}$
is Hausdorff if for all $\lambda\in\Lambda$ we have $X_{\lambda}$ is Hausdorff.
\end{proposition}

\begin{proof}
If $\vec{x}\neq\vec{y}$, then there exists a $\lambda\in\Lambda$ such
that $\pi_{\lambda}\vec{x}\neq\pi_{\lambda}\vec{y}$. Then consider
such a $\lambda$. Therefore since $X_{\lambda}$ is Hausdorff, we can
separate $x_{\lambda}$ and $y_{\lambda}$. Then we have two disjoint
neighborhoods $U\ni x_{\lambda}$ and $V\ni y_{\lambda}$. Then their preimages
$\pi^{-1}_{\lambda}U$ and $\pi^{-1}_{\lambda}V$ are two disjoint
neighborhoods of $\vec{x}$ and $\vec{y}$ (respectively). Hence the
points $\vec{x}$ and $\vec{y}$ are separated. Since they were
arbitrary points, the claim follows.
\end{proof}

\begin{example}
The space $[0,1]^{[0,1]}$ with the product topology is normal $(T_{4})$.
[Fact: compact Hausdorff spaces are normal.] However, $[0,1]^{[0,1]}$
is not metrizable, so it is not first-countable $(A1)$.
\end{example}

\begin{lemma}[Urysohn]
Let $(X,\topology)$ be a normal space. Let $F_{0}$ and $F_{1}$ be two
disjoint subspaces of $X$. Then there exists a continuous function
$f\colon X\to[0,1]$ such that $f|_{F_{0}}=0$ and $f|_{F_{1}}=1$.
\end{lemma}

\begin{theorem}[Tietze]
Let $(X,\topology)$ be a normal space. Let $A\subset X$ be a closed
subset. Let $f\colon(A,\topology_{A})\to[0,1]$ be continuous. Then
there exists a unique continuous function $\overline{f}\colon X\to[0,1]$
such that $\overline{f}|_{A}=f$.
\end{theorem}

(This may be viewed as a generalization of Urysohn's lemma.)

\endinput